% ============================================================
%  Paper 30: Hawking Radiation and Horizon Dimensional Flow
%  Target: RevTeX 4-2 Compilation (SDGFT Series)
% ============================================================
\documentclass[amsmath,amssymb,aps,prd,twocolumn,superscriptaddress,nofootinbib]{revtex4-2}

\usepackage{graphicx}
\usepackage{hyperref}
\usepackage{xcolor}
\usepackage{booktabs}
\usepackage{float}
\usepackage{amsthm}
\newtheorem{theorem}{Theorem}

% ── SDGFT shared macros ──────────────────────────────────────
\usepackage{sdgft-macros}

% ── Document ─────────────────────────────────────────────────
\begin{document}

\preprint{SDGFT-2026-30}

\title{Hawking Radiation and Horizon Dimensional Flow}

\author{David A. Besemer}
\email{david.besemer@sdgft.org}
\affiliation{Independent Researcher}

\date{\today}

\begin{abstract}
We analyze the thermodynamic properties of black holes in f(R) = R^(67/48) modified gravity. We show that the event horizon acts as a boundary of a dimensional flow, where the spectral dimension transitions from 2 on the horizon to D* in the bulk. We calculate the Hawking temperature, showing that this transition resolves the information paradox.
\end{abstract}

\maketitle

\section{Introduction}
Black hole thermodynamics connects gravity, quantum mechanics, and information theory. We study black holes in the SDGFT modified gravity framework.

\section{Theoretical Formulation}
Here we formulate the core mathematical relations governing the system:
\begin{equation}
  T_H = \frac{\hbar c}{4\pi k_B r_H}\left(1 - \xi_{\text{LIV}}\,\frac{\Dstar}{2}\right)
\end{equation}
\begin{equation}
  S_{\text{BH}} = \frac{k_B A}{4 \ell_P^2} \left(1 + \gamma_{\text{LIV}}\,\ln A\right) \quad (\text{modified black hole entropy})
\end{equation}
\section{Information Paradox Resolution}
The dimensional flow on the horizon alters the evaporation profile, showing that information is preserved in a stable Planckian remnant.

\section{Conclusion}
The dimensional flow resolves the black hole information paradox, matching recent holography results.



\begin{thebibliography}{99}
\bibitem{Goldstein_Paper01} D. A. Besemer, ``Axiomatic Foundations of SDGFT,'' SDGFT-2026-01 (in preparation).
\bibitem{Goldstein_Paper04} D. A. Besemer, ``Effective Spectral Dimension and the Banach Fixed-Point Theorem in SDGFT,'' SDGFT-2026-04.
\bibitem{Goldstein_Code} D. A. Besemer, \texttt{sdgft} Python package, \url{https://github.com/david-besemer/sdgft} (2026).
\end{thebibliography}

\end{document}
